% !TeX program = xelatex
% !BIB program = biber
\documentclass[UTF8,a4paper,twocolumn,zihao=5,fontset=none]{ctexart}

% 基于 EmpyreanHYR 双栏中文论文模板的版式思想重构。
% 上游许可与版本见 licenses/upstream_template.md 及
% licenses/upstream_template_LICENSE.txt。

% ---------- 字体：优先课程目标字体，缺失时自动后备 ----------
\usepackage{fontspec}
\usepackage{xeCJK}

\IfFontExistsTF{Times New Roman}{
  \setmainfont{Times New Roman}
}{
  \setmainfont{TeX Gyre Termes}
}

\IfFontExistsTF{SimSun}{
  \setCJKmainfont{SimSun}[BoldFont=SimHei]
}{
  \setCJKmainfont{FandolSong}[BoldFont=FandolSong-Bold]
}
\IfFontExistsTF{SimHei}{
  % SimHei 通常只有常规字形；启用受控伪粗体以避免标题层级调用
  % \bfseries 时发生字形替代警告。
  \setCJKsansfont{SimHei}[AutoFakeBold=2]
  \setCJKfamilyfont{zhhei}{SimHei}[AutoFakeBold=2]
}{
  \setCJKsansfont{FandolHei}[BoldFont=FandolHei-Bold]
  \setCJKfamilyfont{zhhei}{FandolHei}[BoldFont=FandolHei-Bold]
}
\IfFontExistsTF{FandolFang}{
  \setCJKfamilyfont{zhfangsong}{FandolFang}[AutoFakeBold=2]
}{
  \setCJKfamilyfont{zhfangsong}{FandolSong}[BoldFont=FandolSong-Bold]
}
\setCJKmonofont{FandolFang}[AutoFakeBold=2]
\newcommand{\fangsongfont}{\CJKfamily{zhfangsong}}
\newcommand{\heiti}{\CJKfamily{zhhei}}

% ---------- 页面与双栏 ----------
\usepackage[left=2cm,right=2cm,top=2cm,bottom=2cm,
            headheight=14pt,headsep=0.45cm,footskip=0.8cm]{geometry}
\setlength{\columnsep}{0.8cm}
\setlength{\columnseprule}{0pt}
\setlength{\parindent}{2em}
\setlength{\parskip}{0pt}
\linespread{1.08}
\usepackage{indentfirst}
\usepackage{setspace}
\usepackage{ragged2e}
\usepackage{microtype}
\XeTeXlinebreaklocale "zh"
\XeTeXlinebreakskip = 0pt plus 1pt
\widowpenalty=10000
\clubpenalty=10000
\displaywidowpenalty=10000
\raggedbottom

% ---------- 标题层级 ----------
\ctexset{
  section={
    format=\heiti\bfseries\zihao{-4},
    number=\arabic{section},
    aftername=\hspace{0.55em},
    beforeskip=0.85em,
    afterskip=0.35em,
    indent=0pt
  },
  subsection={
    format=\heiti\bfseries\zihao{5},
    number=\thesection.\arabic{subsection},
    aftername=\hspace{0.5em},
    beforeskip=0.65em,
    afterskip=0.25em,
    indent=0pt
  },
  subsubsection={
    format=\heiti\zihao{5},
    number=\thesubsection.\arabic{subsubsection},
    aftername=\hspace{0.45em},
    beforeskip=0.5em,
    afterskip=0.2em,
    indent=0pt
  }
}

% ---------- 数学 ----------
\usepackage{amsmath,amssymb,mathtools,bm}
\numberwithin{equation}{section}

% ---------- 图表 ----------
\usepackage{graphicx}
\usepackage{booktabs,tabularx,array,threeparttable,multirow,makecell}
\usepackage{caption}
\usepackage{subcaption}
\usepackage{tikz}
\usetikzlibrary{arrows.meta,positioning,calc,fit}
\usepackage{siunitx}
\sisetup{detect-all}
\usepackage{needspace}
\usepackage{enumitem}

\renewcommand{\figurename}{图}
\renewcommand{\tablename}{表}
\captionsetup{
  font={small},
  labelfont=bf,
  labelsep=quad,
  justification=centering,
  singlelinecheck=true,
  skip=4pt
}

% ---------- 页眉页脚 ----------
\usepackage{fancyhdr}
\pagestyle{fancy}
\fancyhf{}
\fancyhead[L]{\zihao{-5}\fangsongfont 无线边缘网络中联邦学习通信效率优化方法研究综述}
\fancyhead[R]{\rmfamily\small\thepage}
\renewcommand{\headrulewidth}{0.4pt}
\renewcommand{\footrulewidth}{0pt}

% ---------- GB/T 7714—2015 顺序编码参考文献 ----------
\usepackage[
  backend=biber,
  style=gb7714-2015,
  sorting=none,
  giveninits=true,
  gbpub=false,
  gbnamefmt=lowercase
]{biblatex}
\addbibresource{references.bib}
\setlength{\bibhang}{2em}
\setlength{\bibitemsep}{0.05em}
\renewcommand*{\bibfont}{\zihao{5}\selectfont\setstretch{1.0}}
\AtBeginBibliography{\RaggedRight\setlength{\emergencystretch}{2em}}

% ---------- 超链接与中文交叉引用 ----------
\usepackage[hidelinks,unicode,bookmarks=true,bookmarksnumbered=true]{hyperref}
\usepackage[nameinlink,noabbrev]{cleveref}
\crefname{figure}{图}{图}
\Crefname{figure}{图}{图}
\crefname{table}{表}{表}
\Crefname{table}{表}{表}
\crefname{equation}{式}{式}
\Crefname{equation}{式}{式}
\crefname{section}{第}{第}
\Crefname{section}{第}{第}
\crefname{subsection}{第}{第}
\Crefname{subsection}{第}{第}

% ---------- 常用列类型 ----------
\newcolumntype{Y}{>{\centering\arraybackslash}X}
\newcolumntype{L}{>{\raggedright\arraybackslash}X}


\newcommand{\PaperTitleZh}{无线边缘网络中联邦学习通信效率优化方法研究综述}
\newcommand{\PaperTitleEn}{A Review of Communication-Efficiency Optimization Methods for Federated Learning in Wireless Edge Networks}
\newcommand{\CoFirstMark}{\ensuremath{\uparrow}}
\newcommand{\AuthorZh}{%
  李嘉宇\textsuperscript{1\CoFirstMark}\hspace{2.2em}%
  潘彦哲\textsuperscript{1\CoFirstMark}\hspace{2.2em}%
  张振国\textsuperscript{1\CoFirstMark}\hspace{2.2em}%
  周士栋\textsuperscript{1\CoFirstMark}}
\newcommand{\AuthorPdf}{李嘉宇, 潘彦哲, 张振国, 周士栋}
\newcommand{\AffiliationZh}{\textsuperscript{1}电子科技大学\qquad\textsuperscript{\CoFirstMark}共同第一作者}
\newcommand{\AuthorEn}{%
  Li Jiayu\textsuperscript{1\CoFirstMark},\hspace{0.9em}%
  Pan Yanzhe\textsuperscript{1\CoFirstMark},\hspace{0.9em}%
  Zhang Zhenguo\textsuperscript{1\CoFirstMark},\hspace{0.9em}%
  Zhou Shidong\textsuperscript{1\CoFirstMark}}
\newcommand{\AffiliationEn}{\textsuperscript{1}University of Electronic Science and Technology of China\qquad\textsuperscript{\CoFirstMark}Joint first authors.}

\hypersetup{
  pdftitle={\PaperTitleZh},
  pdfauthor={\AuthorPdf},
  pdfsubject={专业写作基础课程综述论文},
  pdfkeywords={联邦学习, 无线边缘网络, 通信效率, 模型压缩, 空中计算}
}

\begin{document}

% 首页题录与摘要保持通栏，正文从摘要后切换为双栏。
\twocolumn[
  \begin{@twocolumnfalse}
    \begin{center}
      {\heiti\bfseries\zihao{3}\PaperTitleZh\par}
      \vspace{0.65em}
      {\fangsongfont\zihao{5}\AuthorZh\par}
      \vspace{0.25em}
      {\fangsongfont\zihao{5}\AffiliationZh\par}
    \end{center}

    \vspace{0.45em}
    \begin{minipage}{\textwidth}
      \fangsongfont\zihao{5}\setstretch{1.18}
      \noindent\textbf{摘要：}无线边缘网络中的联邦学习以模型交换替代原始数据汇聚，但频谱受限、信道时变和终端异构使多轮通信成为主要瓶颈。本文围绕通信轮数、参与节点、传输载荷、无线传输过程和网络组织方式，梳理客户端选择与频率控制、无线资源分配、模型压缩、空中计算和分层联邦学习五类方法，并比较其收益、代价、假设及适用场景。分析表明，各类方法作用于不同成本分量，通信收益通常以精度损失、控制开销、信道协同或部署复杂度为代价；跨层组合需统一核算比特量、时延、能耗和聚合误差。本文可为无线边缘场景下通信优化方法的选择与组合提供参考。

      \noindent\textbf{关键词：}联邦学习；无线边缘网络；通信效率；空中计算；分层联邦学习
    \end{minipage}

    \vspace{0.45em}
    \begin{center}
      {\rmfamily\bfseries\fontsize{14pt}{17pt}\selectfont\PaperTitleEn\par}
      \vspace{0.3em}
      {\rmfamily\fontsize{10.5pt}{12.5pt}\selectfont\AuthorEn\par}
      {\rmfamily\fontsize{9pt}{11pt}\selectfont\AffiliationEn\par}
    \end{center}

    \vspace{0.25em}
    \begin{minipage}{\textwidth}
      {\rmfamily\fontsize{9pt}{11pt}\selectfont
      \hyphenpenalty=10000\exhyphenpenalty=10000\emergencystretch=1.2em
      \noindent\textbf{Abstract:} Federated learning replaces raw-data collection with iterative model exchange, yet limited spectrum, time-varying channels, and heterogeneous devices make communication a major bottleneck in wireless edge networks. This review organizes existing approaches according to the cost component they primarily affect: communication rounds and participating clients, radio-resource use, transmitted payload, over-the-air aggregation, and network hierarchy. Their benefits, additional costs, operating assumptions, and suitable deployment scenarios are compared under a common framework. The analysis shows that communication gains are commonly traded for learning distortion, control overhead, channel coordination, or deployment complexity. Consequently, cross-layer combinations should be evaluated through joint accounting of transmitted bits, latency, energy consumption, and aggregation error. The review provides a structured basis for selecting and combining communication-efficiency techniques in wireless edge federated learning.

      \noindent\textbf{Keywords:} federated learning; wireless edge networks; communication efficiency; over-the-air computation; hierarchical federated learning}
    \end{minipage}
    \vspace{0.9em}
    \hrule height 0.4pt
    \vspace{0.9em}
  \end{@twocolumnfalse}
]

\section{引言}\label{sec:introduction}

联邦学习将原始数据保留在终端，通过周期性交换模型状态完成协同训练。其通信代价主要来自反复执行的模型下发和更新上传。经典联邦平均（Federated Averaging，FedAvg）在相邻两次聚合之间安排多步本地计算，以终端计算换取较少的全局同步\cite{McMahan2017FedAvg}。

无线边缘网络会放大这一通信负担。参与设备共享有限频谱，链路质量随时间和位置变化；终端的算力、能量、在线状态及本地数据量也不一致。高维模型的多轮交换因而会占用带宽，延长训练等待时间，并消耗终端电量\cite{Lim2020MobileEdge,Niknam2020WirelessCommunications,Zhang2023EdgeFLSurvey}。通信成本可先理解为“轮数、参与节点数与单次载荷”的累积，再由无线速率映射为时延和能耗。因此，通信效率不能只用单轮压缩倍数衡量，还应在相同学习目标下比较累计比特、墙钟时延和通信能耗。

现有综述已经从不同侧面建立了研究基础。移动边缘网络综述讨论通信成本、资源配置和隐私约束\cite{Lim2020MobileEdge}，无线通信综述强调信道与接入资源对联邦训练的影响\cite{Niknam2020WirelessCommunications}。中文研究则分别梳理了边缘智能中的联邦学习、跨设备客户端选择及高效联邦学习方法\cite{Zhang2023EdgeFLSurvey,Zhang2024ClientSelectionSurvey,Ge2025EfficientFLSurvey}。这些工作的覆盖范围和分类尺度不同。若只并列算法名称，读者难以判断方法直接减少的是轮数、节点、载荷、传输时间还是长距离回传。

本文按照方法直接改变的通信成本分量组织文献。综述范围包括通信轮数、参与集合、传输载荷、无线传输过程和网络组织方式，并归纳为客户端选择与频率控制、无线资源分配、模型压缩、空中计算和分层拓扑五类。隐私、安全和激励只在其增加交互或资源约束时讨论。第~\ref{sec:system-model}节建立训练流程和评价口径，第~\ref{sec:methods}节归纳五类方法，第~\ref{sec:comparison}节比较其收益、代价与适用场景，第~\ref{sec:challenges}节讨论共同问题与发展方向，第~\ref{sec:conclusion}节给出结论。

\section{无线边缘网络中的联邦学习通信模型与评价指标}\label{sec:system-model}

\subsection{系统模型与基本训练流程}\label{subsec:training-model}

中心式无线边缘联邦学习由聚合服务器和 \(K\) 个终端组成。服务器可位于接入点、边缘节点或云端；分层系统还设有边缘聚合层\cite{Tran2019WirelessOptimization,Liu2020HierarchicalFL}。一次全局通信轮包含五个步骤：服务器下发当前模型，选择参与客户端；客户端使用本地数据训练若干步并上传更新；服务器聚合这些更新，得到下一轮模型。原始数据留在终端，网络传输的是模型参数、梯度或编码结果\cite{McMahan2017FedAvg}。

为说明各客户端的训练目标如何合成为全局目标，设终端 \(k\) 保存数据集 \(\mathcal D_k=\{\xi_{k,i}\}_{i=1}^{n_k}\)，其中 \(n_k\) 为样本数。全局经验风险写为各客户端局部损失的加权平均：
\begin{equation}\label{eq:global-objective}
  \begin{aligned}
    \min_{\bm w}\quad F(\bm w)
      &= \sum_{k=1}^{K} q_k F_k(\bm w), \\
    F_k(\bm w)
      &= \frac{1}{n_k}\sum_{i=1}^{n_k}
         \ell(\bm w;\xi_{k,i}), \\
    q_k&=\frac{n_k}{N},
    \qquad N=\sum_{j=1}^{K}n_j .
  \end{aligned}
\end{equation}
式中，\(\bm w\) 为模型参数，\(\ell\) 为单样本损失，\(F_k\) 为局部经验风险，\(q_k\) 为样本数权重，\(N\) 为总样本数。用普通语言说，样本较多的客户端对全局目标贡献更大。式~\eqref{eq:global-objective}也给出了通信优化的前提：减少通信后仍要达到约定的损失或精度目标。

设 \(t\) 为全局轮次，\(\mathcal S_t\) 为参与集合，\(e\) 为本地更新步。客户端收到全局模型后，首先令本地初始模型 \(\bm w_{k,t}^{(0)}=\bm w_t\)。随后采用小批量随机梯度下降（stochastic gradient descent，SGD）：每一步只抽取一小批样本估计梯度，再沿降低损失的方向更新参数，
\begin{equation}\label{eq:local-update}
  \begin{aligned}
    \bm w_{k,t}^{(0)}&=\bm w_t,\\
    \bm g_{k,t}^{(e)}
      &=\frac{1}{\lvert\mathcal B_{k,t}^{(e)}\rvert}
        \sum_{\xi\in\mathcal B_{k,t}^{(e)}}
        \nabla\ell\!\left(\bm w_{k,t}^{(e)};\xi\right),\\
    \bm w_{k,t}^{(e+1)}
      &=\bm w_{k,t}^{(e)}-\eta_t\bm g_{k,t}^{(e)}.
  \end{aligned}
\end{equation}
其中，\(\mathcal B_{k,t}^{(e)}\) 是第 \(e\) 步的小批量样本，\(\bm g_{k,t}^{(e)}\) 是估计梯度，\(\eta_t\) 是学习率。式~\eqref{eq:local-update}可以读作：客户端从服务器下发的同一模型出发，用本地小批量梯度逐步更新参数。它说明本地计算发生在两次通信之间，为后文讨论通信频率提供了基础。

客户端完成 \(E_k\) 步本地更新后，联邦平均（Federated Averaging，FedAvg）按参与客户端的数据量计算加权平均\cite{McMahan2017FedAvg}：
\begin{equation}\label{eq:fedavg}
  \bm w_{t+1}
  =\sum_{k\in\mathcal S_t}
   \frac{n_k}{\sum_{j\in\mathcal S_t}n_j}
   \bm w_{k,t}^{(E_k)} .
\end{equation}
式~\eqref{eq:fedavg}可以读作：服务器按参与客户端的样本数加权，得到下一轮全局模型。增加 \(E_k\) 不会取消当前轮的下发或上传，只会增加两次聚合之间的本地计算。若算法仍能收敛，它可能减少达到目标精度所需的全局轮数。非独立同分布（non-IID）指不同客户端的数据分布不相同。此时本地训练过多会使模型沿不同方向移动，形成“本地漂移”，并抵消减少同步的收益\cite{McMahan2017FedAvg,Wang2019AdaptiveFL}。

\subsection{通信成本分解与评价指标}\label{subsec:communication-metrics}

通信效率需要多项指标描述。通信轮数是达到指定精度所需的全局同步次数；参与节点数是每轮上传的客户端数量；单轮比特量包括下发、上传和控制信令。墙钟时延是训练开始至达到目标精度的实际时间，包含计算、传输和同步等待；通信能耗则是无线收发消耗的能量。“发送得少”“发送得快”和“发送得省电”分别对应比特量、时延和能耗，不能相互替代。

为统一核算累计通信量，设 \(T\) 为达到预定精度或损失阈值所需的全局轮数，\(B_t^{\mathrm{dl}}\) 为第 \(t\) 轮下行广播量，\(B_{k,t}^{\mathrm{ul}}\) 为客户端 \(k\) 的上行载荷，\(B_{k,t}^{\mathrm{ctl}}\) 为调度、索引和信道反馈等控制比特，则
\begin{equation}\label{eq:total-bits}
  \begin{split}
    B_{\mathrm{tot}}=\sum_{t=0}^{T-1}\biggl[
      B_t^{\mathrm{dl}}+\sum_{k\in\mathcal S_t}
      \bigl(B_{k,t}^{\mathrm{ul}}+B_{k,t}^{\mathrm{ctl}}\bigr)
    \biggr].
  \end{split}
\end{equation}
式~\eqref{eq:total-bits}可以读作：逐轮累加一次下行广播，以及所有参与客户端的上传和控制信息。\textbf{说明性示例：}若训练进行 2 轮，每轮下行广播 120 bit，3 个客户端各上传 100 bit 更新和 10 bit 控制信息，则累计通信量为 \(2\times[120+3\times(100+10)]=900\) bit。该计算只用于解释公式，不是实验结果。它说明单轮载荷变小不保证 \(B_{\mathrm{tot}}\) 下降；所需轮数或控制信息增加时，累计通信量仍可能上升\cite{Ge2025EfficientFLSurvey}。

五类方法作用于式~\eqref{eq:total-bits}的不同部分。客户端选择与频率控制改变 \(\mathcal S_t\) 和 \(T\)；模型压缩改变载荷，也可能增加索引等控制比特；无线资源分配通常不改变载荷，而是改变传输时间和能量。空中计算让多用户传输与聚合同时进行，更适合用信道占用时间和聚合误差衡量\cite{Zhang2023EdgeFLSurvey}。分层联邦学习改变消息路径，需要分别统计端—边和边—云通信。

无线链路把比特量进一步映射为时延和能耗。在一个信道状态近似不变的上行时隙内，令 \(b_{k,t}\)（\si{\hertz}）为分配带宽，\(P_{k,t}\)（\si{\watt}）为发射功率，\(h_{k,t}\) 为信道功率增益，\(N_0\)（\si{\watt\per\hertz}）为噪声功率谱密度，\(I_{k,t}\)（\si{\watt}）为干扰功率，则常用关系为\cite{Tran2019WirelessOptimization,Dinh2021WirelessConvergence}
\begin{subequations}\label{eq:wireless-link}
  \begin{align}
    R_{k,t}&=b_{k,t}\log_2\!\left(
      1+\frac{P_{k,t}h_{k,t}}
      {N_0b_{k,t}+I_{k,t}}\right),\label{eq:uplink-rate}\\
    \tau_{k,t}^{\mathrm{ul}}
      &=\frac{B_{k,t}^{\mathrm{ul}}}{R_{k,t}},\label{eq:uplink-delay}\\
    E_{k,t}^{\mathrm{ul}}
      &=P_{k,t}\tau_{k,t}^{\mathrm{ul}}.\label{eq:uplink-energy}
  \end{align}
\end{subequations}
式~\eqref{eq:wireless-link}中，\(R_{k,t}\) 是上行速率，\(\tau_{k,t}^{\mathrm{ul}}\) 是传输时延，\(E_{k,t}^{\mathrm{ul}}\) 是发送能耗。三式可以依次读作：带宽、功率和信道条件决定速率；载荷除以速率得到发送时间；功率乘发送时间得到能耗。提高功率可能加快速率，却也提高单位时间功耗，所以总能耗未必下降。下行广播的比特量通常按基站一次发送统计，而广播时延可能由最慢接收客户端决定。前者统计发送多少，后者描述等待多久。后文讨论资源分配时将沿用这一区分\cite{Dinh2021WirelessConvergence,Zhang2023EdgeFLSurvey}。

\section{通信效率优化方法}\label{sec:methods}

由式~\eqref{eq:total-bits}和式~\eqref{eq:wireless-link}可知，通信开销可以从参与节点、交互频率、单次载荷、无线传输过程和消息路径五个层面处理。下文以方法直接改变的变量为分类依据，而不按算法名称逐篇介绍。

\subsection{客户端选择与通信频率控制}\label{subsec:selection-frequency}

移动边缘训练中，可用设备常多于基站能够同时接入的数量，部分设备还会离线或逾期。客户端选择决定参与集合 \(\mathcal S_t\)，频率控制决定本地步数 \(E_k\)、触发条件和更新时序。两者分别回答“选谁”和“何时通信”，影响参与规模、全局轮数或等待时间。

客户端选择分为随机、资源感知和数据效用感知三类。随机选择无需详细状态，但可能选中慢设备。资源感知选择根据算力、电量、在线状态和信道质量，筛选能按期完成训练与上传的设备。数据效用感知还估计样本覆盖、局部损失或更新价值\cite{Nishio2019FedCS,Yang2020Scheduling,Xu2021ClientBandwidth,Zhang2024ClientSelectionSurvey}。“快设备优先”虽能缩短等待，却会产生选择偏差：设备性能和位置可能与数据分布相关，长期排除弱设备会使参与数据偏离总体。

通信时序可分为同步、异步和事件触发机制。同步训练等待本轮更新到齐后再聚合，轮次时间受最慢参与者限制。异步训练允许更新到达后立即或分批采用，主要减少同步空等。更新陈旧度是客户端所用模型版本与服务器当前版本的差距，差距越大，更新越可能偏离当前方向。异步不必然减少累计上传；若每次更新仍单独发送，它减少的是等待而非比特。事件触发只在模型变化或梯度价值超过阈值时上传，才可能过滤冗余通信\cite{Lu2020EAFLM,Liu2021AsyncFL}。

相关研究已从“让本轮按期完成”扩展到长期公平和动态到达。FedCS按资源状态选择设备，后续无线调度与长期选择再考虑信道、带宽和跨轮参与；自适应方法则按资源预算调整聚合周期\cite{Nishio2019FedCS,Yang2020Scheduling,Xu2021ClientBandwidth,McMahan2017FedAvg,Wang2019AdaptiveFL}。异步边缘学习通过阈值或部分聚合抑制低价值更新，通信—计算重叠与公平调度用于减少等待并保留弱设备机会\cite{Lu2020EAFLM,Liu2021AsyncFL,Ozfatura2021FedTOE}。这些机制分别处理逾期参与、同步空等或冗余上传，并不同时降低所有指标。

这类方法需要状态上报和调度计算，还可能产生陈旧更新与公平性问题。随机选择适合设备条件相近的系统；资源感知选择适合期限严格的场景；数据效用感知适合数据差异明显且效用可估计的任务。异步机制适合设备间歇在线，但不适合对陈旧更新敏感的训练。参与集合确定后，设备仍要共享频谱，下一步还需配置无线与计算资源。

\subsection{无线资源分配与联合优化}\label{subsec:resource-allocation}

摄像头、车辆或手机共同训练时，弱信道或低算力终端仍会拖慢整轮训练。同步聚合需等待更新到齐，最慢终端的“本地计算时间加上传时间”常决定轮次完成时间。资源分配配置带宽 \(b_{k,t}\)、功率 \(P_{k,t}\)、接入时隙、中央处理器（CPU）频率和边缘关联，以改善传输时延、能耗或掉队率。

这些变量解决的问题不同。带宽决定频谱份额，功率影响接收质量与速率，时隙决定发送时间；三者作用于无线传输。CPU频率改变本地计算速度和计算能耗，不直接进入上行速率式。边缘关联决定终端连接哪个基站或边缘服务器，进而影响链路质量、服务器负载和辅助计算位置。系统通常先收集信道、能量和计算状态，再在时延或能量约束下求解配置\cite{Tran2019WirelessOptimization,Dinh2021WirelessConvergence}。这些操作通常不减少 \(B_{k,t}^{\mathrm{ul}}\)，主要改变传完相同比特所需的时间和能量。

资源决策分为单轮配置和跨轮配置。单轮方法依据当前状态控制本轮时延、能耗或掉队数量，适合状态变化较慢且要求快速求解的场景。跨轮方法还考虑剩余电量、历史参与和未来信道，避免持续透支同一批设备。相关研究由此扩展到收敛约束、信道调度和长期客户端—带宽设计\cite{Tran2019WirelessOptimization,Dinh2021WirelessConvergence,Yang2020Scheduling,Xu2021ClientBandwidth}。

多个联邦任务共享接入网时，资源配置还要在业务之间分配频谱。边缘辅助方案会同时决定弱终端是否由边缘侧辅助，并配置CPU频率、无线资源和模型稀疏度\cite{Xu2022MultiFLBandwidth,Chen2025EdgeAssistedSparsification}。这类联合决策能同时处理通信等待、计算等待和长期能量约束，但变量增多会提高在线求解难度。在线学习或强化学习虽能适应动态状态，也会产生探索、训练和控制开销。

资源优化适合蜂窝边缘、多业务接入和设备差异明显的系统。若网络稳定、带宽充足或状态难以获取，复杂优化的收益可能不足以抵消开销。信道反馈占用控制比特，在线求解消耗时间和算力，求解期间状态变化还可能使配置过期。因此，反馈和调度应计入 \(B_{k,t}^{\mathrm{ctl}}\)，比特、时延和能耗也应分别报告。资源分配解决“怎样把更新传完”；若更新本身过大，还需用压缩削减载荷。

\subsection{模型压缩与通信载荷削减}\label{subsec:model-compression}

深度模型参数很多，即使信道较好，完整上传模型或梯度仍会占用大量带宽。模型压缩直接改变单轮载荷：量化减少每个数值使用的比特数，稀疏化减少需要发送的元素数量。量化先把连续数值映射到有限等级，再发送等级编码；稀疏化先按幅值或阈值保留少量重要元素，其余元素本轮不发送。两者可以组合，但会分别引入量化误差和丢弃误差。

为区分数学稀疏性与真实传输开销，设待传更新为 \(\bm g_{k,t}\in\mathbb R^d\)，稀疏算子为 \(S(\cdot)\)，本文定义
\begin{equation}\label{eq:effective-compression}
  \begin{aligned}
    \sigma_{k,t}
      &=1-\frac{\lVert S(\bm g_{k,t})\rVert_0}{d},\\
    \rho_{k,t}^{\mathrm{eff}}
      &=\frac{B_{k,t}^{\mathrm{val}}+B_{k,t}^{\mathrm{idx}}
      +B_{k,t}^{\mathrm{hdr}}}{B_{k,t}^{\mathrm{raw}}}.
  \end{aligned}
\end{equation}
式~\eqref{eq:effective-compression}第一行中，\(d\) 是原始元素总数，零范数 \(\lVert\cdot\rVert_0\) 统计保留的非零元素数，\(\sigma_{k,t}\) 表示被省略元素的\mbox{比例}。第二行比较实际编码量与未压缩载荷：\(B^{\mathrm{val}}\) 是保留数值的比特，\(B^{\mathrm{idx}}\) 是位置索引，\(B^{\mathrm{hdr}}\) 是码长、形状等头部信息，\(B^{\mathrm{raw}}\) 是原始载荷。有效传输比 \(\rho_{k,t}^{\mathrm{eff}}\) 越小，真实节省越多。稀疏化必须发送索引，否则服务器无法知道每个数值对应哪个参数；因此理论稀疏率高，不等于真实传输比一定同比降低。

误差反馈用于补偿暂未发送的信息。客户端把上一轮残差加入当前更新，再压缩上传；原更新与发送近似值之差留在本地，下一轮继续参与计算。被舍弃的小更新因而能在后续补偿。客户端为此需要保存残差并增加筛选或编码计算；终端频繁退出时，残差状态也难以连续维护\cite{Sattler2020STC,Han2020AdaptiveSparsification,Tian2022GradientCompression}。

代表性工作体现了从单一压缩到自适应压缩的演进。FedPAQ把周期平均、部分参与和量化放入同一框架；一比特压缩感知继续降低数值位宽，但需要稀疏先验和服务器重构\cite{Reisizadeh2020FedPAQ,Zhang2022OneBitCS}。稀疏三值压缩结合双向稀疏、低比特表示和误差累积，自适应方法则按训练状态调整保留比例；近期工作还同时选择稀疏度、量化级和能耗目标\cite{Sattler2020STC,Han2020AdaptiveSparsification,Tian2022GradientCompression,Marnissi2024AdaptiveSQ,Chen2025EdgeAssistedSparsification}。这些方法直接减少式~\eqref{eq:total-bits}中的载荷，而不是加快既定比特的传输。

压缩适合高维模型、窄上行且终端能够承担编码的场景。对小模型、短消息或弱算力设备，索引、头部和编解码成本可能抵消收益。压缩过强会增加优化误差，压缩感知或安全聚合还会引入重构、码本及额外交互\cite{Chen2022SecureCompressedFL}。评价时应同时报告累计比特、精度或损失、编解码成本，并说明是否计入索引。压缩回答“每个客户端发送多少”，空中计算则改变“多个客户端如何同时发送并聚合”。

\subsection{空中计算与通信聚合融合}\label{subsec:aircomp}

数字正交上传需要服务器分别恢复每个客户端的更新，因而通常为终端分配不同的时隙或子载波，避免数字码字相互干扰。服务器完成解码后再进行加权，而FedAvg等模型聚合本质上就是对客户端更新求加权和。无线多址信道又具有信号叠加特性：若终端把更新映射为模拟符号并同时发送，接收波形便包含这些符号的和。空中计算（over-the-air computation，AirComp）据此把“分别传输后求和”改为“传输过程中求和”\cite{Amiri2019OTA,Zhu2020BroadbandAnalog,ZhuLi2020IntegratedAirComp}。两种机制如图~\ref{fig:digital-aircomp}所示。

\begin{figure*}[t]
  \centering
  \resizebox{0.94\textwidth}{!}{\begin{tikzpicture}[
  font=\small,
  >=Latex,
  panel/.style={draw=black!55, rounded corners=2pt, minimum width=7.9cm, minimum height=3.55cm},
  client/.style={draw, rounded corners=1.5pt, minimum width=1.45cm, minimum height=0.52cm, align=center, fill=black!3},
  block/.style={draw, rounded corners=1.5pt, minimum width=1.65cm, minimum height=0.58cm, align=center, fill=black!7},
  sum/.style={draw, circle, minimum size=0.7cm, inner sep=0pt, fill=black!7},
  flow/.style={->, line width=0.55pt}
]
  \node[panel] (leftpanel) at (-4.15,0) {};
  \node[panel] (rightpanel) at (4.15,0) {};
  \node[font=\bfseries] at (-4.15,1.43) {数字正交传输};
  \node[font=\bfseries] at (4.15,1.43) {空中模拟聚合};

  \node[client] (ld1) at (-7.15,0.72) {终端 1};
  \node[client] (ld2) at (-7.15,0) {终端 2};
  \node[client] (ld3) at (-7.15,-0.72) {终端 \(K\)};
  \node[block] (dec1) at (-4.55,0.72) {逐用户解码};
  \node[block] (dec2) at (-4.55,0) {逐用户解码};
  \node[block] (dec3) at (-4.55,-0.72) {逐用户解码};
  \node[sum] (lsum) at (-1.55,0) {\(\sum\)};
  \draw[flow] (ld1) -- node[above, font=\scriptsize] {资源 1} (dec1);
  \draw[flow] (ld2) -- node[above, font=\scriptsize] {资源 2} (dec2);
  \draw[flow] (ld3) -- node[above, font=\scriptsize] {资源 \(K\)} (dec3);
  \draw[flow] (dec1) -- (lsum);
  \draw[flow] (dec2) -- (lsum);
  \draw[flow] (dec3) -- (lsum);
  \draw[flow] (lsum) -- ++(0.72,0) node[right, font=\scriptsize] {聚合值};
  \node[font=\scriptsize, align=center] at (-4.15,-1.34) {独立时隙或子载波：先可靠解码，再在应用层求和};

  \node[client] (ra1) at (1.12,0.72) {终端 1};
  \node[client] (ra2) at (1.12,0) {终端 2};
  \node[client] (ra3) at (1.12,-0.72) {终端 \(K\)};
  \node[sum] (channel) at (4.18,0) {\(\sum\)};
  \node[block] (scale) at (6.25,0) {接收缩放\(a_t\)};
  \draw[flow] (ra1) -- node[above, sloped, font=\scriptsize] {\(h_{1,t}\beta_{1,t}x_{1,t}\)} (channel);
  \draw[flow] (ra2) -- node[above, font=\scriptsize] {\(h_{2,t}\beta_{2,t}x_{2,t}\)} (channel);
  \draw[flow] (ra3) -- node[below, sloped, font=\scriptsize] {\(h_{K,t}\beta_{K,t}x_{K,t}\)} (channel);
  \draw[flow] (channel) -- node[above, font=\scriptsize] {\(y_t\)} (scale);
  \draw[flow] (scale) -- ++(0.82,0) node[right, font=\scriptsize] {\(\widehat{s}_t\)};
  \draw[flow] (4.18,1.0) -- node[right, font=\scriptsize] {噪声 \(n_t\)} (channel);
  \node[font=\scriptsize, align=center] at (4.15,-1.34) {同频同时发送：同步、信道估计与功率对齐后在信号域求和};
\end{tikzpicture}
}
  \caption{数字正交上传与空中聚合的机理差异\protect\cite{Amiri2019OTA,Zhu2020BroadbandAnalog,ZhuLi2020IntegratedAirComp}}
  \label{fig:digital-aircomp}
\end{figure*}

为说明信道如何影响加权和，下面采用经过相位补偿的等效实数信道模型。对第 \(t\) 轮某一模型坐标，参与集合为 \(\mathcal S_t\)，终端 \(k\) 发送归一化更新 \(x_{k,t}\)，目标聚合权重为 \(q_{k,t}\)，则
\begin{subequations}\label{eq:aircomp}
  \begin{align}
    y_t &= \sum_{k\in\mathcal S_t}h_{k,t}\beta_{k,t}x_{k,t}+n_t,\label{eq:aircomp-receive}\\
    \widehat{s}_t &= a_t y_t,\qquad
    s_t=\sum_{k\in\mathcal S_t}q_{k,t}x_{k,t},\label{eq:aircomp-target}\\
    e_t &= \widehat{s}_t-s_t
    =\sum_{k\in\mathcal S_t}
      \bigl(a_t h_{k,t}\beta_{k,t}-q_{k,t}\bigr)x_{k,t}+a_t n_t.\label{eq:aircomp-error}
  \end{align}
\end{subequations}
式~\eqref{eq:aircomp-receive}表示接收端看到所有终端信号与噪声的叠加。信道系数 \(h_{k,t}\) 描述幅度变化，终端预编码 \(\beta_{k,t}\) 用于发送前补偿，\(n_t\) 是接收噪声。式~\eqref{eq:aircomp-target}中，接收缩放 \(a_t\) 调整整体幅度，\(\widehat{s}_t\) 是实际估计，\(s_t\) 是目标加权和。式~\eqref{eq:aircomp-error}进一步说明：聚合误差由系数未对齐和缩放后的噪声组成，理想对齐要求 \(a_t h_{k,t}\beta_{k,t}\approx q_{k,t}\)。这正是后文功率控制、设备选择和接收机设计需要减小的误差。

数字方案需要分离多用户信号，AirComp则把受控叠加直接用作计算结果。它的收益来自无线资源复用，而不是无损压缩更新比特。

对齐并非总能实现。信道反演令预编码系数与 \(h_{k,t}\) 近似成反比，但深衰落终端可能需要超过其最大发射功率。此时只能截断预编码、降低公共幅度或暂停该终端。信道估计不准会使 \(a_t h_{k,t}\beta_{k,t}\) 偏离目标权重；载波频率、相位或符号时刻不同步，也会破坏同一坐标上的叠加。接收噪声则直接进入式~\eqref{eq:aircomp-error}。这些因素形成的是聚合误差，不能简单依靠应用层重传消除。

AirComp适合终端多、聚合函数可加、允许小幅模拟误差，且基站能获得较准信道信息并维持同步的场景。宽带及时变信道研究通过功率控制、设备选择和接收缩放降低聚合均方误差；稀疏一比特和分层AirComp还结合压缩或网络组织，但会叠加量化、信道与层级误差\cite{Tegin2023TimeVaryingOTA,Oh2024SparseOneBitOTA,Azimi2024HierarchicalOTA}。

若深衰落频繁、同步难以保证，或任务要求可靠恢复每个客户端更新，数字上传的独立解码与差错控制更稳妥。AirComp主要减少多用户接入时间，不能理解为无条件、零误差或零能耗的聚合。

\subsection{分层联邦学习与网络拓扑优化}\label{subsec:hierarchical-fl}

终端直接连接中心云时，模型更新要跨越接入网和回传网。远距离链路增加时延、流量和掉线风险，云端还要等待分散终端。分层联邦学习把客户端按小区、位置或边缘关联划分为簇，由邻近边缘服务器先聚合。多数高频更新因而停留在端—边短链路，只有簇模型以较低频率经过边—云回传\cite{Liu2020HierarchicalFL,Abad2020HierarchicalCellular}。“端—边高频、边—云低频”表示一个云端轮次内进行多次簇内聚合，随后才同步云端，如图~\ref{fig:hierarchical-aggregation}所示。

\begin{figure*}[t]
  \centering
  \resizebox{0.62\textwidth}{!}{\begin{tikzpicture}[
  font=\small,
  >=Latex,
  cloud/.style={draw, rounded corners=2pt, minimum width=2.15cm, minimum height=0.62cm, align=center, fill=black!10},
  edge/.style={draw, rounded corners=1.5pt, minimum width=1.85cm, minimum height=0.58cm, align=center, fill=black!6},
  client/.style={draw, rounded corners=1.2pt, minimum width=1.25cm, minimum height=0.5cm, align=center, fill=black!2},
  cluster/.style={draw=black!55, dashed, rounded corners=2pt, inner sep=5pt},
  flow/.style={->, line width=0.55pt}
]
  \node[cloud] (cloud) at (0,3.15) {云端服务器\\全局聚合};
  \node[edge] (e1) at (-2.25,1.65) {边缘服务器 1};
  \node[edge] (e2) at (2.25,1.65) {边缘服务器 \(M\)};
  \node[client] (c11) at (-3.1,0.25) {终端 1};
  \node[client] (c12) at (-1.4,0.25) {终端 \(K_1\)};
  \node[client] (c21) at (1.4,0.25) {终端 1};
  \node[client] (c22) at (3.1,0.25) {终端 \(K_M\)};

  \node[cluster, fit=(e1)(c11)(c12), label={[font=\scriptsize]below:边缘簇 1}] (g1) {};
  \node[cluster, fit=(e2)(c21)(c22), label={[font=\scriptsize]below:边缘簇 \(M\)}] (g2) {};

  \draw[flow] (c11) -- (e1);
  \draw[flow] (c12) -- (e1);
  \draw[flow] (c21) -- (e2);
  \draw[flow] (c22) -- (e2);
  \draw[flow] (e1) -- node[above left, font=\scriptsize, align=center] {边—云同步\\每 \(H\) 轮一次} (cloud);
  \draw[flow] (e2) -- node[above right, font=\scriptsize, align=center] {边—云同步\\每 \(H\) 轮一次} (cloud);
  \node[font=\scriptsize, align=center] at (0,0.92) {端—边交换：高频本地更新与簇内聚合};
\end{tikzpicture}
}
  \caption{端—边—云双时间尺度聚合\protect\cite{Liu2020HierarchicalFL,Abad2020HierarchicalCellular,Liu2023HierarchicalQuantization}}
  \label{fig:hierarchical-aggregation}
\end{figure*}

为区分三种时间尺度，令 \(t\) 表示云端轮次，\(r=1,\ldots,H\) 表示其中的边缘聚合序号，\(e=0,\ldots,E_{k,m,t}^{(r)}\) 表示本地SGD步数。这里 \(H\) 只表示每个云端轮次包含的边缘聚合次数，不表示本地步数。设系统有 \(M\) 个边缘簇，第 \(m\) 个簇的客户端集合为 \(\mathcal C_{m,t}\)，本地轨迹为 \(\bm u_{k,m,t}^{(r,e)}\)，最终本地模型记为 \(\bar{\bm u}_{k,m,t}^{(r)}=\bm u_{k,m,t}^{(r,E_{k,m,t}^{(r)})}\)。两级聚合可写为
\begin{subequations}\label{eq:hierarchical-aggregation}
  \begin{align}
    \bm w_{m,t}^{(0)}&=\bm w_t,\qquad
    \bm u_{k,m,t}^{(r,0)}=\bm w_{m,t}^{(r-1)},\label{eq:hierarchical-init}\\
    N_{m,t}&=\sum_{k\in\mathcal C_{m,t}}n_k,\qquad
    \bm w_{m,t}^{(r)}
    =\sum_{k\in\mathcal C_{m,t}}\frac{n_k}{N_{m,t}}
      \bar{\bm u}_{k,m,t}^{(r)},\label{eq:edge-aggregation}\\
    N_t&=\sum_{m=1}^{M}N_{m,t},\qquad
    \bm w_{t+1}
    =\sum_{m=1}^{M}\frac{N_{m,t}}{N_t}\bm w_{m,t}^{(H)}.\label{eq:cloud-aggregation}
  \end{align}
\end{subequations}
式~\eqref{eq:hierarchical-init}表示云模型先下发到边缘，并成为客户端的初始模型。式~\eqref{eq:edge-aggregation}表示边缘服务器按样本数汇总簇内模型。式~\eqref{eq:cloud-aggregation}表示完成 \(H\) 次边缘聚合后，云端再把各簇模型合成为下一轮全局模型。三式依次对应图~\ref{fig:hierarchical-aggregation}中的云端下发、端—边训练和边—云同步。

增大 \(H\) 会在两次云端同步之间执行更多簇内聚合，因此减少边—云同步频率。各簇也会沿自身数据分布独立演化更久；簇间数据差异较大时，边缘模型可能向不同方向移动，形成簇间漂移，并需要更多云端轮次重新对齐。客户端分簇会影响数据代表性，边缘负载不均会产生新的慢节点，终端移动还会改变集合 \(\mathcal C_{m,t}\) 并增加模型迁移与控制通信。

分层方法适合多小区蜂窝、工业物联网和云—边—端集群，其收益取决于回传瓶颈、簇内数据分布和边缘可用性。量化分层可同时选择聚合间隔和消息精度；异构物联网方案还要处理异步到达、边缘负载和客户端选择\cite{Liu2023HierarchicalQuantization,Tang2023AdaptiveGradient,Wang2023HierarchicalIoT}。端—边链路也可结合压缩或AirComp，并用资源分配平衡各簇\cite{Azimi2024HierarchicalOTA}。组合收益不能简单相加，需要分别核算两级比特、时延与能耗，并报告压缩误差、信道失真和层级漂移。

\section{方法比较与适用场景分析}\label{sec:comparison}

\subsection{统一维度下的横向比较}\label{subsec:unified-comparison}

横向比较需区分直接改变的变量与最终系统指标。图~\ref{fig:communication-cost-map}把式~\eqref{eq:total-bits}中的轮数、参与集合和上传载荷分别对应频率控制、客户端选择和压缩；速率与时延、并发聚合及消息路径则对应资源分配、AirComp和分层联邦学习。

\begin{figure*}[t]
  \centering
  \resizebox{0.96\textwidth}{!}{\begin{tikzpicture}[
  font=\small,
  >=Latex,
  root/.style={draw, rounded corners=2pt, minimum width=5.2cm,
               minimum height=0.78cm, align=center, fill=black!9,
               font=\bfseries},
  item/.style={draw, rounded corners=1.5pt, text width=2.25cm,
               minimum height=1.15cm, align=center, fill=black!3,
               inner sep=3pt},
  flow/.style={->, line width=0.55pt}
]
  \node[root] (root) at (0,1.75) {联邦训练的通信成本与完成代价};
  \coordinate (bus) at (0,0.95);
  \draw[line width=0.55pt] (root.south) -- (bus);
  \draw[line width=0.55pt] (-6.45,0.95) -- (6.45,0.95);

  \node[item] (round) at (-6.45,0) {全局轮数 \(T\)\\[2pt]\textbf{通信频率控制}};
  \node[item] (client) at (-3.87,0) {参与集合 \(\mathcal S_t\)\\[2pt]\textbf{客户端选择}};
  \node[item] (payload) at (-1.29,0) {单次载荷 \(B_{k,t}^{\mathrm{ul}}\)\\[2pt]\textbf{量化与稀疏化}};
  \node[item] (radio) at (1.29,0) {速率 \(R\)、时延、能耗\\[2pt]\textbf{无线资源分配}};
  \node[item] (aircomp) at (3.87,0) {同时接入与聚合误差\\[2pt]\textbf{AirComp}};
  \node[item] (hier) at (6.45,0) {消息路径与回传次数\\[2pt]\textbf{分层联邦学习}};

  \foreach \n in {round,client,payload,radio,aircomp,hier}{
    \draw[flow] (\n |- bus) -- (\n.north);
  }
\end{tikzpicture}
}
  \caption{通信成本分量与五类优化方法的对应关系\protect\cite{Lim2020MobileEdge,Zhang2023EdgeFLSurvey,Ge2025EfficientFLSurvey}}
  \label{fig:communication-cost-map}
\end{figure*}

同一方法可能影响多个指标，但分类应以其直接控制的变量为准：压缩改变比特，资源分配改变速率和等待。表~\ref{tab:method-comparison}据此按“变量—原理—收益—代价”比较五类方法。

\begin{table*}[t]
  \centering
  \caption{五类通信效率优化方法的统一比较}
  \label{tab:method-comparison}
  \begin{threeparttable}
    \fontsize{7.1pt}{8.6pt}\selectfont
    \setlength{\tabcolsep}{2.8pt}
    \renewcommand{\arraystretch}{1.22}
    \begin{tabularx}{\textwidth}{
      @{}>{\centering\arraybackslash}p{1.35cm}
      >{\centering\arraybackslash}p{1.25cm}
      >{\raggedright\arraybackslash}p{1.72cm}
      >{\raggedright\arraybackslash}p{1.62cm}
      >{\raggedright\arraybackslash}p{1.62cm}
      >{\raggedright\arraybackslash}p{1.68cm}
      >{\raggedright\arraybackslash}p{1.62cm}
      >{\raggedright\arraybackslash}X@{}}
      \toprule
      方法类别 & 优化对象 & 代表机制 & 主要收益 & 额外代价 & 关键假设 & 适用场景 & 主要局限 \\
      \midrule
      客户端选择与频率控制
        & \(T,\mathcal S_t,E_k\)
        & 效用选择；异步或事件触发
        & 少上传、少同步、少等待
        & 状态上报、偏差与陈旧
        & 状态可观测；更新具代表性
        & 大规模异构物联网
        & 弱设备可能长期缺席，影响公平与收敛\cite{Nishio2019FedCS,Wang2019AdaptiveFL,Zhang2024ClientSelectionSurvey} \\
      无线资源分配
        & 带宽、功率、时隙、CPU频率及边缘关联
        & 跨轮资源与关联联合调度
        & 降低时延、能耗和掉队率
        & 信道反馈、求解和信令
        & 状态可估计；链路模型有效
        & 蜂窝、多业务、无人机网络
        & 改变传输效率但未必减少累计比特\cite{Dinh2021WirelessConvergence,Xu2022MultiFLBandwidth,Fan2025EdgeAssociation,Meng2025UAVEnergy} \\
      模型与梯度压缩
        & \(B_{k,t}^{\mathrm{ul}}\)、稀疏度和量化位宽
        & 量化、稀疏、误差反馈
        & 减少载荷和发送能耗
        & 失真、索引和重构计算
        & 更新可压缩；误差可补偿
        & 高维模型、窄带终端
        & 理论压缩率不等于有效传输比\cite{Reisizadeh2020FedPAQ,Sattler2020STC,Marnissi2024AdaptiveSQ} \\
      AirComp
        & 多用户上传时间、频谱及聚合误差
        & 同步模拟发送；波形域求和
        & 传输与聚合并行
        & 信道状态、同步、功控、噪声
        & 函数可加；物理层可协同
        & 密集同步无线接入
        & 深衰落和时变信道会破坏对齐\cite{Zhu2020BroadbandAnalog,Tegin2023TimeVaryingOTA,Oh2024SparseOneBitOTA} \\
      分层联邦学习
        & 消息路径、簇划分及两级聚合间隔 \(H\)
        & 边缘高频、云端低频聚合
        & 少回传、易扩展
        & 部署、分组、控制和漂移
        & 边缘可用；簇权重合理
        & 多小区、云—边—端集群
        & 簇偏差、节点迁移和故障增加协调难度\cite{Liu2020HierarchicalFL,Liu2023HierarchicalQuantization,Wang2023HierarchicalIoT,Zhong2022HierarchicalHeterogeneous} \\
      \bottomrule
    \end{tabularx}
    \begin{tablenotes}[flushleft]
      \footnotesize
      \item 注：\(T\) 为全局轮数，\(\mathcal S_t\) 为参与集合，\(E_k\) 为本地步数；表中内容不用于比较不同实验设置中的百分比。
    \end{tablenotes}
  \end{threeparttable}
\end{table*}


表~\ref{tab:method-comparison}表明，通信收益通常伴随成本转移：少参与带来偏差风险，压缩和AirComp引入失真，资源与分层优化增加控制开销。比特减少不必然缩短端到端时延，单轮加速也不保证总能耗下降\cite{Tran2019WirelessOptimization,Ge2025EfficientFLSurvey}，因而不能脱离目标精度和无线基线判断最优类别。

\subsection{场景匹配与方法组合}\label{subsec:scenario-matching}

方法选择应先识别主导瓶颈，再决定是否组合多个优化层。表~\ref{tab:scenario-selection}中的“优先”只表示先处理主导成本。若设备可用性还受参与意愿影响，激励应作为客户端选择的约束，而非新的通信类别\cite{Jia2023HierarchicalIncentive}。

\begin{table*}[t]
  \centering
  \caption{不同无线边缘场景下的方法选择}
  \label{tab:scenario-selection}
  \begin{threeparttable}
    \small
    \setlength{\tabcolsep}{4pt}
    \renewcommand{\arraystretch}{1.28}
    \begin{tabularx}{\textwidth}{
      @{}>{\centering\arraybackslash}p{2.05cm}
      >{\raggedright\arraybackslash}p{2.45cm}
      >{\centering\arraybackslash}p{2.65cm}
      >{\raggedright\arraybackslash}X
      >{\raggedright\arraybackslash}X@{}}
      \toprule
      主导瓶颈 & 典型场景 & 优先考虑的方法 & 选择理由 & 使用时需要注意的问题 \\
      \midrule
      同步等待与终端异构
        & 大规模物联网；设备能力和在线状态差异明显
        & 客户端选择＋资源分配
        & 同时控制“谁参与”和“如何完成本轮传输”，缓解慢终端造成的等待
        & 快设备优先会造成选择偏差；状态采集和联合求解也有开销\cite{Xu2021ClientBandwidth,Nishio2019FedCS} \\
      高维更新与窄上行
        & 视觉模型、移动终端或低速上行链路
        & 量化、稀疏化及误差反馈
        & 直接减少每轮需要编码和发送的更新量
        & 需计入索引、码本和残差；压缩误差可能增加训练轮数\cite{Reisizadeh2020FedPAQ,Sattler2020STC} \\
      密集终端同时接入
        & 同步条件可控、聚合函数可加的蜂窝或传感网络
        & AirComp；必要时结合稀疏化
        & 在同一时频资源中完成多用户叠加，避免逐用户正交接入
        & 依赖信道估计、同步与功率可行性，并引入聚合误差\cite{Zhu2020BroadbandAnalog,Oh2024SparseOneBitOTA} \\
      回传链路受限
        & 多小区、工业园区或云—边—端集群
        & 分层联邦＋两级压缩
        & 先减少边—云同步频率，再降低簇模型的回传载荷
        & 需联合控制量化误差、簇间漂移和边缘负载\cite{Liu2023HierarchicalQuantization,Zhong2022HierarchicalHeterogeneous} \\
      多瓶颈随时间变化
        & 移动、多业务竞争或时变信道环境
        & 跨轮选择、压缩和资源联合优化
        & 可随训练阶段和网络状态切换主要优化变量
        & 在线决策、反馈和模型迁移可能抵消部分通信收益\cite{Xu2021ClientBandwidth,Marnissi2024AdaptiveSQ,Azimi2024HierarchicalOTA} \\
      \bottomrule
    \end{tabularx}
  \end{threeparttable}
\end{table*}


客户端选择与无线资源分配的组合同时回答“选谁”和“如何传”。前者筛除本轮难以完成训练的终端，后者为入选设备配置带宽、功率和时隙；但两层若都偏向快设备，会放大数据代表性偏差，状态上报与在线求解也会形成控制流量\cite{Xu2021ClientBandwidth,Nishio2019FedCS}。

稀疏化与AirComp同时降低有效维度和多用户接入时间。终端需协调稀疏位置或采用可聚合编码，否则不同支持集无法按坐标叠加；量化误差、残差近似、信道失配和噪声还会共同进入聚合误差\cite{Oh2024SparseOneBitOTA}。

分层联邦与压缩先以低频云端同步减少回传次数，再压缩簇模型以降低单次回传载荷。两级压缩误差会与簇间漂移叠加，边缘侧还要承担编解码、簇管理和负载均衡\cite{Liu2023HierarchicalQuantization,Chen2025EdgeAssistedSparsification}。因此，应分别核算端—边和边—云开销，不能直接相加不同基线下的收益。

\section{现存问题与发展趋势}\label{sec:challenges}

\subsection{现有研究的共同局限}\label{subsec:common-limitations}

指标口径不统一会限制横向比较。轮数、压缩率、时延、能耗和聚合误差衡量的对象不同；若目标精度以及控制消息、索引和重传的统计范围不同，相同“通信降低”并不表示相同系统收益\cite{Lim2020MobileEdge,Mills2020WirelessEdgeIoT,Ge2025EfficientFLSurvey}。

理想信道假设和终端异构都会限制部署结论。固定拓扑、准确信道状态和同步终端便于分析，但移动、深衰落和故障会改变参与集合、AirComp对齐与分层归属\cite{Tegin2023TimeVaryingOTA,Azimi2024HierarchicalOTA}。设备能力还可能与数据分布相关，快设备优先、长期本地训练和固定分簇都可能造成统计偏差。

跨层组合会叠加误差和控制开销：压缩失真、信道噪声、更新陈旧与簇间漂移可能同时出现，状态反馈和在线求解也会消耗资源。区块链异步方案还需执行模型上链、共识和信誉评估，产生额外通信与计算\cite{Huang2024BlockchainAFL}。秘密共享和安全传输同样只作为效率约束讨论，不构成独立主线\cite{Chen2022SecureCompressedFL,Meng2025UAVEnergy}。

\subsection{面向可部署系统的发展方向}\label{subsec:future-directions}

发展方向应回应上述限制。在线跨层方法需要让客户端选择、压缩程度、无线资源和边缘关联随训练阶段与信道变化，同时计入决策时间和反馈比特\cite{Xu2021ClientBandwidth,Marnissi2024AdaptiveSQ,Azimi2024HierarchicalOTA}。面向异构终端的方法还要同时处理设备能力、模型规模和数据分布差异。边缘辅助与集群蒸馏依赖额外边缘算力或代理知识，其收益仍需跨场景验证\cite{Zhong2022HierarchicalHeterogeneous,Chen2025EdgeAssistedSparsification,Zhang2026FedCD}。

面向任务的消息设计不再要求每轮传输完整参数。单轮通信方案让客户端上传生成模型，由服务器生成模拟数据并训练全局模型\cite{Zheng2023SingleRound}；集群蒸馏则交换任务知识\cite{Zhang2026FedCD}。这类证据仍受具体任务和生成质量限制。可复现实验应在相同目标精度下报告累计比特、墙钟时延、能耗、控制开销和聚合误差，并公开信道轨迹、数据划分与重复设置。

\section{结论}\label{sec:conclusion}

通信效率方法可归纳为客户端选择与频率控制、无线资源分配、模型压缩、AirComp和分层拓扑五类。它们分别作用于轮数与参与节点、传输载荷、无线过程和消息路径，也会引入统计偏差、失真、控制信令或部署复杂度。窄带高维模型可优先考虑压缩；密集且同步条件可控的接入场景可考虑AirComp；多小区或回传受限系统更适合分层聚合。组合方法应先确定主导瓶颈，再在同一目标精度下分别核算比特、时延、能耗和误差，不能把不同实验基线下的收益直接相加。


\FloatBarrier
\printbibliography[heading=bibintoc,title={参考文献}]

\end{document}
